\ifx\fontfam\undefined
\else
  \fontfam[LM]
  \def\sl{\slant\it}
\fi

\enlang \uselanguage{ukenglish}
\hyperlinks{}{}

\parindent=0pt

{\bf HUGO HOUSE YEARLONG WRITE THE FIRST DRAFT OF YOUR NOVEL, 2026/2027 BOOK PROPOSAL}\par 
\bigskip

{\bf PAUL M. CRAY}\par 
\bigskip

{\bf GENRE:}\par
Historiographic metafiction (the novel is metafictional in that it is about itself as a novel, and historiographic in that it is about the technoeconomic and sociocultural forces that drive history). Alternatively, though the novel contains both science fiction (low magic) and fantasy (high magic), it is a fabulation, defined by John Clute in {\sl The Encyclopedia of Science Fiction} as ‘any story which challenges the two main assumptions of genre sf: that the world can be seen; and that it can be told.’\par
\bigskip

{\bf TITLE:}\par
{\sl The Visibilities}\par
\bigskip

{\bf LOGLINE:}\par
The travails of one man, and his avatars scattered across the multiverse, in attempting to overcome his cosmic loneliness and awake from the nightmare of history.\par
\bigskip

{\bf COMP TITLES:}\par
Henry Adams, ‘A Law of Acceleration’ and ‘The Rule of Phase Applied to History’\par
Kingsley Amis, {\sl Lucky Jim}\par
Martin Amis, {\sl Money}\par
Isaac Asimov, {\sl Extraterrestrial Civilisations}\par
Margaret Atwood, {\sl The Handmaid's Tale}\par
Julian Barnes, {\sl A History of the World in 10½ Chapters}\par
Jorge Luis Borges, {\sl Labyrinths}\par
William Cooper, {\sl Scenes from Provincial Life}\par
Len Deighton, {\sl Bomber}\par
BBC, {\sl Doctor Who}\par
A. S. Byatt, {\sl Possession}\par
John Clute and David Langford (eds), {\sl The Encyclopedia of Science Fiction}\par
John Clute and John Grant (eds), {\sl The Encyclopedia of Fantasy}\par
Jonathan Coe, {\sl What a Carve Up!} (US title: {\sl The Winshaw Legacy: or, What a Carve Up!)}\par
Umberto Eco, {\sl The Name of the Rose}\par
Richard P. Feynman, Robert B. Leighton and Matthew Sands, {\sl The Feynman Lectures on Physics}\par
C.B.P. Finn, {\sl Thermal Physics}\par
Ian Fleming, {\sl From Russia, with Love}\par
Michael Frayn, {\sl Sweet Dreams}\par
Alan Garner, {\sl Red Shift}\par
Alasdair Gray, {\sl Lanark: A Life in Four Books}\par
Garth Risk Hallberg, {\sl City on Fire}\par
Douglas R. Hofstadter, {\sl Gödel, Escher, Bach: An Eternal Golden Braid}\par
B.S. Johnson, {\sl Travelling People}\par
Uwe Johnson, {\sl Anniversaries}\par
James Joyce, {\sl Ulysses}\par
Donald E. Knuth, {\sl The \TeX book}\par
Stanley Kubrick and Arthur C. Clarke, {\sl 2001: A Space Odyssey}\par
Ursula K. Le Guin, {\sl A Wizard of Earthsea}\par
David Lodge, {\sl How Far Can You Go?} (US title: {\sl Souls and Bodies})\par
Marc W. Miller, {\sl Traveller: Science-Fiction Adventure in the Far Future}\par
David Mitchell, {\sl Cloud Atlas}\par
Christopher Priest, {\sl The Prestige}\par
Thomas Pynchon, {\sl Gravity's Rainbow}\par
Frank Richards, {\sl Billy Bunter and the Greyfriars Pretender}\par
Olaf Stapledon, {\sl Star Maker}\par
Laurence Sterne, {\sl The Life and Opinions of Tristram Shandy, Gentleman}\par
Neal Stephenson, {\sl Cryptonomicon}\par
Nicholas Tomalin and Ron Hall, {\sl The Strange Last Voyage of Donald Crowhurst}\par
Jenny Turner, {\sl London Review of Books} articles (\url{https://www.lrb.co.uk/contributors/jenny-turner})\par
Vernor Vinge, {\sl A Fire Upon the Deep}\par
Kurt Vonnegut Jr., {\sl Slaughterhouse-Five, or The Children's Crusade: A Duty-Dance with Death}\par
David Foster Wallace, {\sl Infinite Jest}\par
Evelyn Waugh, {\sl Sword of Honour} trilogy\par
{\sl Wisden Cricketers' Almanack 1977}, 114th edition\par
\bigskip

{\bf BACK COVER COPY:}\par
The main text of {\sl The Visibilities} is forced into a 366-element tensor that becomes so overloaded that inevitably it bursts out of its confines to infect the paratext of the novel, including its peritext. {\sl The Visibilities} is both a grimoire and a teratography. The front cover is a representation of a region of the Heptabrot set procedurally generated with METAPOST. At the centre is the hexagonal ‘fort’ surrounded by a beautifully rich, colourful, whirling, roiling sevenfold structure. The hexagonal fort is in the centre of the cover, the AR (Atomic Razor) device physically impressed (blind blocked) into the board (of the hardback, like the AP device on old Academic Press books). The back cover is similar, but shows a region of the Octabrot set. Here there is an eightfold structure surrounding the central—for this Multibrot set—heptagonal ‘fort’. Ideally, matters will be arranged so that the sevenfold and eightfold structures seamlessly meld into one another, but the realisation of that element of the book is left for the future. The front and spine will carry the name of the author and title (plus the publisher's device). There will be no text on the back cover. The book will be sold new shrink-wrapped so that the barcode will go on the disposable outer covering, so that it is not corrupted by any extraneous commercial clutter.\par
\bigskip

\bye
